% Research notes for integration into the ProveIt polylogarithms manuscript.
% Preamble requirements: amsmath, amssymb, amsthm; theorem/proposition/corollary.
% Scope: independent derivations. Anderson/Ouyang antecedents explicitly credited.
\section{Integral distribution coordinates and primitive-grid torsion}
\label{intdist:sec:main}

The complex character coordinates of the manuscript diagonalize the
weighted distribution relations, including their exceptional weights.
They do not by themselves determine the integral lattice or justify
reduction of that diagonalization modulo a prime dividing a character
group order. We give a concrete integral model and then compute a
primitive-grid cokernel completely at a product of two distinct primes.
The latter calculation reveals torsion invisible to complex rank and
jet computations.

The general freeness and resolution method belongs to the established
theory of universal distributions, originating with Kubert's ordinary
distribution theorem \cite{intdist:Kubert1979}. In particular, Ouyang's universal
norm distribution has a free-basis theorem and an Anderson resolution
\cite[Proposition 4.2 and Theorem 5.1]{intdist:Ouyang2002}; the ordinary
case is treated in \cite{intdist:Ouyang2001}. The purpose of the first
part below is to supply an explicit, finite, all-divisor grid version
with arbitrary weights, arbitrary coefficient rings, and actual grid
points as basis vectors. The two-prime cokernel, its integral Smith
factors, and its sharp denominator consequence are the specific
additional calculations obtained here. No global priority assertion is
made for these formulas.

\subsection{A basis which survives every specialization}

Let $q=\prod_{p\mid q}p^{e_p}\ge2$, let $R$ be any commutative unital
ring, and choose $t_p\in R$. Put $G_q=(1/q)\mathbb Z/\mathbb Z$,
$V_q=R[G_q]$, and
\[
 D_p(x)=\sum_{\substack{y\in G_q\\py=x}}[y]-t_p[x],
 \qquad x\in G_{q/p},\qquad
 U_q(R,\mathbf t)=V_q/\langle D_p(x)\rangle.
\]
Thus $G_{q/p}=pG_q$; no distribution row outside its legitimate source
is used. All tensor products in this section are over $R$.

For $m=p^e$, define the following set of local residues:
\begin{equation}\label{intdist:eq:local-basis}
 B(p^e)=\{0\}\cup\{p^{e-1}+1,p^{e-1}+2,\ldots,p^e-1\}.
\end{equation}
An empty upper interval is allowed; in particular $B(2)=\{0\}$.
Use the additive primary decomposition of $G_q$ to put
\begin{equation}\label{intdist:eq:grid-basis}
 \mathcal B_q=
 \left\{\left[\sum_{p\mid q}\frac{b_p}{p^{e_p}}\right]:
                b_p\in B(p^{e_p})\right\}.
\end{equation}
Distinct choices give distinct grid points.

\begin{theorem}[Integral basis and arbitrary base change]
\label{intdist:thm:basis}
For every $R$ and every choice of the weights, the images of
$\mathcal B_q$ form an $R$-basis of $U_q(R,\mathbf t)$. Consequently
\[
 U_q(R,\mathbf t)\cong R^{\varphi(q)}.
\]
This basis is independent of the weights and contains $[0]$. In
particular, imposing the homogeneous endpoint anchor $[0]=0$ leaves a
free module of rank $\varphi(q)-1$ over every $R$.

Over the universal ring $\mathbb Z[t_p:p\mid q]$, every grid point has
unique integral polynomial coordinates in this basis. For a point
$x=\sum_px_p$ whose $p$-primary component has additive order $p^{j_p}$,
each coordinate polynomial has degree at most $j_p$ in $t_p$.
\end{theorem}

For Hurwitz evaluation, use the representative of a grid point in
$(0,1]$, with $[0]$ evaluated as $\zeta(s,1)=\zeta(s)$.
The term ``integral'' refers to formal coordinates. It asserts neither
linear independence of evaluated zeta or polylogarithm values nor an
integral-basis assertion for a field generated by such values.

\subsubsection{The finite resolution and its filtration}

Order the prime divisors of $q$. For $S\subseteq\{p:p\mid q\}$ set
$d_S=\prod_{p\in S}p$ and
\[
 C_r=\bigoplus_{|S|=r}R[G_{q/d_S}].
\]
Write a generator of its $S$-summand as $[x,S]$. For $p\in S$, define
\[
 \partial_p[x,S]=
 \sum_{\substack{y\in G_{q/d_{S\setminus\{p\}}}\\py=x}}
          [y,S\setminus\{p\}]
 -t_p[x,S\setminus\{p\}],
\]
and let $\partial$ be their sum with the usual exterior-algebra signs.
The transfers and inclusions for distinct primes commute. For example,
composing two transfers enumerates exactly the $p\ell$ preimages of
$x$, while a transfer by $p$ cannot raise the $\ell$-primary order when
$p\ne\ell$. Thus $\partial^2=0$ and
$H_0(C_\bullet)=U_q(R,\mathbf t)$.

For a $p$-primary point $x_p$, let $h_p(x_p)$ be its additive-order
exponent, with $h_p(0)=0$. Give $[x,S]$ filtration degree
\begin{equation}\label{intdist:eq:height}
 |S|+\sum_{p\mid q}h_p(x_p).
\end{equation}
The degree lies between zero and $\sum_pe_p$. In a $p$-transfer, every
preimage other than the zero preimage of $x_p=0$ raises $h_p$ by one;
all the other primary heights are unchanged. The leading part of
$\partial$ preserves \eqref{intdist:eq:height}. The inclusion term and
the zero preimage lower it by one.

There is a minor coordinate twist which must be accounted for: a
$p$-preimage multiplies every other primary component by $p^{-1}$.
In the $S$-summand replace its $\ell$-primary coordinate by
\[
 x_\ell\longmapsto
 \left(\prod_{p\in S,\ p\ne\ell}p\right)^{-1}x_\ell.
\]
All these inverses are units on the relevant primary cyclic group, so
this is merely a permutation of the grid basis. It preserves heights.
In the new coordinates, the leading $p$-transfer acts only on the
$p$-primary factor. In degree zero, where $S$ is empty, this coordinate
change is the identity.

The associated graded complex is therefore the tensor product, over
$p\mid q$, of the two-term complexes
\begin{equation}\label{intdist:eq:local-complex}
 R[G_{p^{e_p-1}}]\xrightarrow{N_p^+}R[G_{p^{e_p}}],
 \qquad
 N_p^+[x]=
 \sum_{\substack{py=x\\h_p(y)=h_p(x)+1}}[y].
\end{equation}
The local domain is in homological degree one, with its height shifted
by one. The supports of distinct $N_p^+[x]$ are disjoint. Each support
contains exactly one residue in $\{1,\ldots,p^{e_p-1}\}$; this gives a
coefficient-one pivot. Hence $N_p^+$ is split injective, and its free
cokernel has the actual local grid basis \eqref{intdist:eq:local-basis}.

A tensor product of these split two-term complexes has zero positive
homology and degree-zero homology freely generated by
$\mathcal B_q$. This can be seen without a flatness hypothesis: split
each local target into the image of $N_p^+$ and its chosen free
complement. The tensor product becomes a direct sum of its free
degree-zero complement and contractible complexes.

For completeness, finite-filtration lifting gives the requisite
conclusion over any ring. If $z$ is a positive-degree cycle, its leading
filtration term is a boundary in the associated graded complex. Lift a
preimage and subtract its actual boundary from $z$. This lowers the
filtration, so induction makes $z$ an actual boundary. Similarly, if
$\partial y\in C_0$ has lower filtration than $y\in C_1$, the leading
term of $y$ is a degree-one cycle for the leading differential. Subtract
the boundary of a degree-two lift, leaving $\partial y$ unchanged and
lowering the filtration of $y$. Repetition proves that the image of
$\partial:C_1\to C_0$ is strict for the filtration. Therefore the
associated graded of $U_q$ is precisely the degree-zero homology
computed above. Its specified basis lifts to an actual basis: leading
terms prove independence, and descending filtration proves spanning.
There are
$\prod_p(p^{e_p}-p^{e_p-1})=\varphi(q)$ basis points. This proves the
basis assertion and also proves the following resolution statement.

\begin{corollary}[All relation syzygies]
\label{intdist:cor:resolution}
The augmented finite complex
\[
 0\longrightarrow C_{\omega(q)}\longrightarrow\cdots
 \longrightarrow C_1\longrightarrow C_0
 \longrightarrow U_q(R,\mathbf t)\longrightarrow0
\]
is exact. It is split exact as a complex of $R$-modules. In particular,
every relation among prime-distribution rows is generated by the
commuting-prime two-step relations, and all higher relations are
accounted for by the same complex.
\end{corollary}

To justify the last splitting assertion, $U_q$ is free, so
$C_0\to U_q$ splits. Its kernel is projective; the surjection from
$C_1$ onto that kernel then splits. Continue inductively. No claim about
semisimplicity or character projectors is involved.

\subsubsection{A terminating polynomial reduction algorithm}

Represent a point by a tuple $x=(a_p\bmod p^{e_p})_p$. If every local
residue belongs to $B(p^{e_p})$, it is already a basis point. Otherwise
choose any local pivot $1\le a_p\le p^{e_p-1}$. All the other
$p$-preimages of $p x$ have local $p$-residues outside this pivot
interval. The distribution row rewrites
\begin{equation}\label{intdist:eq:rewrite}
 [x]=t_p[p x]-
 \sum_{j=1}^{p-1}
 \left[x+\frac{j}{p}\right],
\end{equation}
where every point is taken in $G_q$ and the $p$-primary component of
$x$ is the chosen pivot. The right-hand siblings may include the point
with zero $p$-component. That point has lower primary height.

The nonzero siblings keep the total primary height and remove one
pivot. The term $p x$ lowers the $p$-primary height by exactly one;
multiplication by $p$ preserves all the other primary heights.
Consequently the lexicographically ordered pair
\[
 (\hbox{total primary height},\ \hbox{number of local pivots})
\]
strictly decreases in every recursive branch. Reduction terminates,
and uniqueness follows from Theorem~\ref{intdist:thm:basis}. Every
occurrence of $t_p$ lowers the $p$-height by one. This proves the degree
bound asserted in that theorem. A memoized implementation only needs
the $q$ grid states and performs exact additions and monomial shifts.

\begin{example}[A uniform four-point coordinate grid at level twelve]
With the primes ordered as $2,3$, the basis is
\[
 [0],\quad[2/3],\quad[3/4],\quad[5/12].
\]
Write $a=t_2$, $b=t_3$. For example, reduction gives
\begin{align}
 [1/12]&=a(a-1)b[0]-(b+1)[3/4]-[5/12],\label{intdist:eq:q12-one}\\
 [1/6]&=a(b-1)[0]-(a+1)[2/3],\\
 [7/12]&=a(b-a)[0]-a(a+1)[2/3]+(b+1)[3/4]+[5/12],\\
 [11/12]&=-a(b-1)[0]+a(a+1)[2/3]-[5/12].
\end{align}
These are identities over $\mathbb Z[a,b]$, including $a=b=1$ and
characteristic two. Analytic evaluation substitutes $t_p=p^s$ for
Hurwitz zeta and $t_p=p^{1-s}$ for cyclotomic polylogarithms. At a
spectral pole, evaluate the full meromorphic identity before taking a
limit; a vanishing coefficient times $\zeta(s)$ can have a nonzero
finite contribution.
\end{example}

\subsection{The integral meaning of the primitive determinant}

Let $W_q=R^{U(q)}$ be the formal span of the primitive $q$-grid points,
and let $P_q:W_q\to U_q$ be their coordinate map in the integral basis
above. This coordinate map need not be injective after a singular
specialization. Over $\mathbb Z[t_p]$, it is a square matrix of integral
polynomials. The complex Fourier/conductor calculation in the
manuscript gives
\[
 \Delta_q=\prod_{p^e\Vert q}
 t_p^{\varphi(q/p^e)(p^{e-1}-1)}
 (t_p^{h_p}-1)^{\varphi(q/p^e)/h_p},
 \qquad h_p=\operatorname{ord}_{q/p^e}(p),
\]
where the order modulo one is one. We now remove the ambiguity from
its choice of complex bases.

\begin{theorem}[Integral primitive determinant]
\label{intdist:thm:determinant}
With any fixed orderings of the integral and primitive grid bases,
\begin{equation}\label{intdist:eq:determinant}
 \det P_q=\pm\Delta_q.
\end{equation}
Consequently, for integer weights with $\Delta_q\ne0$, the integral
index of the primitive-grid span is exactly $|\Delta_q|$. For an
arbitrary coefficient ring the primitive grid is a basis if and only
if $\Delta_q$ is a unit.
\end{theorem}

\begin{proof}
Over $\mathbb C[t_p]$, both the integral basis established above and
the conductor basis of the manuscript are bases of the same free
module. Their transition determinant is a nonzero constant, since the
only units of this polynomial ring are nonzero constants. The
primitive Fourier transform is also constant. The manuscript's
diagonal determinant therefore implies
$\det P_q=c\Delta_q$ for a nonzero constant $c$.
The highest monomial of $\Delta_q$ has coefficient one, whereas
$\det P_q$ has integer coefficients. Hence $c\in\mathbb Z$.

It remains to show that no prime divides $c$. Over any field $K$,
regard the weights as algebraically independent. Put $z_p=t_p^{-1}$
and, for $u\in U(q)$, form the formal multivariate series
\[
 F_{r,u}(\mathbf z)=
 \sum_{\substack{(n_p)\in\mathbb Z_{\ge0}^{\omega(q)}\\
    u\prod_pp^{n_p}\equiv r\pmod q}}
 \prod_pz_p^{n_p}.
\]
The primitive rows form the identity matrix. If $p\mid r$, every term
has $n_p\ge1$, so removing one factor of $p$ proves exactly the
$p$-distribution equation. These are rational functions: on the finite
residue grid, the pushforward $M_p$ satisfies
$M_p^{e_p+h_p}=M_p^{e_p}$, and summing its powers gives a rational
resolvent with denominator $1-z_p^{h_p}$. Thus this construction is
valid over $K(t_p)$ in every characteristic and gives an extension of
arbitrary primitive-grid data. Since the full quotient has dimension
$\varphi(q)$ there, $P_q$ is invertible over $K(t_p)$.

If a prime divided $c$, reducing $\det P_q=c\Delta_q$ modulo that
prime would give the zero polynomial, contradicting this generic
invertibility. Therefore $c=\pm1$. The remaining assertions follow
from the determinant criterion for a square matrix over a ring and,
over $\mathbb Z$, the index formula for a full-rank sublattice.
\end{proof}

\subsection{Complete two-prime primitive cokernels}

Assume now that $q=p\ell$ for distinct primes. Write
\begin{equation}\label{intdist:eq:two-prime-notation}
 \begin{gathered}
 a=t_p,\quad b=t_\ell,\qquad
 h_p=\operatorname{ord}_{p}(\ell),\quad
 h_\ell=\operatorname{ord}_{\ell}(p),\\
 c_p=(p-1)/h_p,\qquad c_\ell=(\ell-1)/h_\ell,\\
 S_p=1+b+\cdots+b^{h_p-1},\quad
 S_\ell=1+a+\cdots+a^{h_\ell-1},\\
 D_p=b^{h_p}-1=(b-1)S_p,\qquad
 D_\ell=a^{h_\ell}-1=(a-1)S_\ell.
 \end{gathered}
\end{equation}
Let $T_q=\operatorname{coker}P_q$.

\begin{theorem}[Two-prime cokernel over an arbitrary ring]
\label{intdist:thm:two-prime}
There is an $R$-module isomorphism
\begin{equation}\label{intdist:eq:cokernel}
 T_{p\ell}\cong
 (R/(D_p))^{c_p-1}\oplus(R/(D_\ell))^{c_\ell-1}
 \oplus C(a,b),
\end{equation}
where the three-generator module $C(a,b)$ has generators $e,X,Y$ and
relations
\begin{equation}\label{intdist:eq:core}
 (a-1)(b-1)e=0,\qquad
 S_pX=(a-1)e,\qquad
 S_\ell Y=(b-1)e.
\end{equation}
No weight, geometric sum, or integer group order is inverted.
\end{theorem}

\begin{proof}
Kill the primitive $p\ell$-grid symbols in the original presentation.
The remaining generators are $e=[0]$, the $p-1$ points of exact
denominator $p$, and the $\ell-1$ points of exact denominator $\ell$.
The source of a $p$-distribution row has denominator either one or
$\ell$, and the source of an $\ell$-row has denominator either one or
$p$. Thus every relation is one of
\begin{align*}
 \sum_{u\in U(p)}[u/p]&=(a-1)e,&
 \sum_{v\in U(\ell)}[v/\ell]&=(b-1)e,\\
 [\ell^{-1}u/p]&=b[u/p],&
 [p^{-1}v/\ell]&=a[v/\ell].
\end{align*}
In the last two rows the inverses are residue-class inverses.

Multiplication by $\ell^{-1}$ on $U(p)$ has $c_p$ cycles, each of
length $h_p$. Along one cycle the relations successively eliminate
all but a single generator $X_i$, with the closing relation
$D_pX_i=0$. Its contribution to the first norm row is $S_pX_i$.
The other layer similarly gives generators $Y_j$, relations
$D_\ell Y_j=0$, and norm contribution $S_\ell Y_j$.

The integer coordinate change from $X_1,\ldots,X_{c_p}$ to
$X=\sum_iX_i$ and any $c_p-1$ of the original generators is unimodular.
Since every cycle has the same annihilator, the $c_p-1$ other
coordinates split off as $R/(D_p)$. Apply the same operation to the
$Y$-cycles. The remaining presentation has the two norm relations in
\eqref{intdist:eq:core}, together with $D_pX=0$ and $D_\ell Y=0$.
Modulo the norm relations, these last two are both equivalent to
$(a-1)(b-1)e=0$, because
\[
 D_pX=(b-1)S_pX=(a-1)(b-1)e,
 \quad
 D_\ell Y=(a-1)S_\ell Y=(a-1)(b-1)e.
\]
The replacement works in both directions and uses no cancellation.
This proves the asserted presentation over arbitrary $R$.
\end{proof}

\subsubsection{The exact ordinary-weight torsion and modular ranks}

\begin{corollary}[Ordinary-weight primitive torsion]
\label{intdist:cor:ordinary}
For $R=\mathbb Z$ and $a=b=1$,
\begin{equation}\label{intdist:eq:ordinary}
 T_{p\ell}\cong
 \mathbb Z^{c_p+c_\ell-1}
 \oplus\mathbb Z/h_p\mathbb Z
 \oplus\mathbb Z/h_\ell\mathbb Z.
\end{equation}
A cyclic factor of order one means the zero module. For a rational
prime $r$, the primitive-grid inclusion over $\mathbb F_r$ at these
weights has rank
\begin{equation}\label{intdist:eq:modular-rank}
 \varphi(p\ell)-(c_p+c_\ell-1)
 -\mathbf1_{r\mid h_p}-\mathbf1_{r\mid h_\ell}.
\end{equation}
\end{corollary}

\begin{proof}
At $a=b=1$ the direct cyclic factors in
\eqref{intdist:eq:cokernel} are free copies of $\mathbb Z$, and the
core relations are simply $h_pX=0$, $h_\ell Y=0$, with $e$ free.
This gives \eqref{intdist:eq:ordinary}. Tensor its presentation with
$\mathbb F_r$. Each free coordinate contributes one dimension to the
cokernel, and a finite cyclic coordinate contributes one exactly when
$r$ divides its order. Subtract from the full quotient dimension
$\varphi(p\ell)$ furnished by Theorem~\ref{intdist:thm:basis}.
\end{proof}

For example, $q=15$ has $h_3=2$, $h_5=4$, and $c_3=c_5=1$.
Its ordinary primitive cokernel is
$\mathbb Z\oplus\mathbb Z/2\mathbb Z\oplus\mathbb Z/4\mathbb Z$.
The primitive rank is seven in characteristic zero and five in
characteristic two. At $q=21$, the corresponding cokernel is
$\mathbb Z^2\oplus\mathbb Z/6\mathbb Z$, giving ranks ten in
characteristic zero and nine in characteristics two and three.
For $q=35$, it is
$\mathbb Z\oplus\mathbb Z/4\mathbb Z\oplus\mathbb Z/6\mathbb Z$.
These are failures of a particular integral coordinate inclusion; the
full weighted distribution quotient remains free in every case.

\subsubsection{Smith factors at nonexceptional integer weights}

\begin{corollary}[Explicit Smith factors for the core]
\label{intdist:cor:smith}
Let $a,b\in\mathbb Z$, and assume $D_pD_\ell\ne0$. Define positive
integers
\[
 g_1=\gcd(a-1,b-1,S_p,S_\ell),\qquad
 g_2=\gcd(D_p,D_\ell,S_pS_\ell).
\]
Then
\begin{equation}\label{intdist:eq:smith-core}
 C(a,b)\cong
 \mathbb Z/g_1\mathbb Z\oplus
 \mathbb Z/(g_2/g_1)\mathbb Z\oplus
 \mathbb Z/(|D_pD_\ell|/g_2)\mathbb Z.
\end{equation}
Together with \eqref{intdist:eq:cokernel}, this specifies the entire
finite primitive-grid cokernel as a direct sum of cyclic groups.
Its order is
$|D_p|^{c_p}|D_\ell|^{c_\ell}$ and its exponent is
\begin{equation}\label{intdist:eq:exponent}
 \boxed{\displaystyle
 E_{p\ell}(a,b)=
 \operatorname{lcm}\bigl(|D_p|,|D_\ell|,|(a-1)(b-1)|\bigr).}
\end{equation}
\end{corollary}

\begin{proof}
A row-relation matrix for the core is
\[
 \begin{pmatrix}
 (a-1)(b-1)&0&0\\
 -(a-1)&S_p&0\\
 -(b-1)&0&S_\ell
 \end{pmatrix}.
\]
The greatest common divisor of its entries is $g_1$. The gcd of its
$2\times2$ minors is $g_2$: the three potentially necessary minors
are $D_\ell$, $D_p$, and $S_pS_\ell$, and every other minor is zero or
a multiple of one of these. Its determinant is $D_pD_\ell$.
The determinantal-divisor characterization of integer Smith form
therefore gives \eqref{intdist:eq:smith-core}.
The last invariant factor is a multiple of both $D_p$ and $D_\ell$,
since $g_2$ divides both, so the extra cyclic summands in
\eqref{intdist:eq:cokernel} do not change the exponent. Finally,
\[
 \frac{|D_pD_\ell|}{\gcd(D_p,D_\ell,S_pS_\ell)}
 =\operatorname{lcm}\bigl(|D_p|,|D_\ell|,|(a-1)(b-1)|\bigr),
\]
by taking prime valuations, or by using
$D_pD_\ell=(a-1)(b-1)S_pS_\ell$.
\end{proof}

\subsection{Sharp rational reduction denominators}

Suppose the integer weights are nonexceptional and express every grid
point in the primitive basis over $\mathbb Q$. Let $E$ be the least
positive integer clearing all entries of the resulting extension
matrix. Because $\mathcal B_q$ itself consists of actual grid points,
this is also the least positive integer for which
$E P_q^{-1}$ is integral. Equivalently it is the exponent of
$\operatorname{coker}P_q$.

\begin{corollary}[Exact two-prime denominator]
\label{intdist:cor:denominator}
At $q=p\ell$ and integer weights $a,b$ with $D_pD_\ell\ne0$, the
least common denominator of all primitive-grid extension coefficients
is precisely \eqref{intdist:eq:exponent}. In particular, for Hurwitz
weights $a=p^k$, $b=\ell^k$, $k\ge1$, it is
\begin{equation}\label{intdist:eq:hurwitz-denominator}
 \operatorname{lcm}\left(
 \ell^{k\operatorname{ord}_p(\ell)}-1,
 p^{k\operatorname{ord}_\ell(p)}-1,
 (p^k-1)(\ell^k-1)\right).
\end{equation}
\end{corollary}

The sharpness also follows directly from the manuscript's finite
residue formula. Points of exact denominator $p$ have nonzero
coefficients with reduced denominator $D_p$, points of exact
denominator $\ell$ give $D_\ell$, and the endpoint coefficient is
$1/((a-1)(b-1))$. The relevant one-prime numerators are powers of the
corresponding weight, which are coprime to its power minus one.
Thus each of the three integers in the least common multiple is
required. The Smith calculation determines substantially more than
this denominator alone.

At $q=15$, $k=1$, the exact denominator is
$\operatorname{lcm}(24,80,8)=240$, whereas the general product bound
in the manuscript is $24\cdot80=1920$. The full primitive index is
1920 and its nontrivial Smith factors are $2,4,240$.
At $q=6$, $k=5$, the exact denominator is
$\operatorname{lcm}(242,1023,7502)=22506$, whereas the product bound
is $247566$. Here the Smith factors are $11,22506$.

\subsubsection{A nonsingular level-fifteen Hurwitz example}

The universal integral basis at level fifteen is, in the chosen order,
\[
 [0],\ [2/5],\ [3/5],\ [4/5],\ [2/3],\ [1/15],\ [4/15],\ [7/15].
\]

\begin{corollary}[Explicit reductions at spectral order two]
\label{intdist:cor:q15s2}
Write $Z_r=\zeta(2,r/15)$ for $1\le r<15$, and $Z_0=\zeta(2)$.
Then
\begin{align}
\zeta(2,1/5)
&=\frac{729(Z_1+Z_{11})+81(Z_2+Z_7)
                +9(Z_4+Z_{14})+(Z_8+Z_{13})}{6560},
\label{intdist:eq:q15-fifth}\\
\zeta(2,1/3)
&=\frac{25(Z_1+Z_4+Z_7+Z_{13})
                  +(Z_2+Z_8+Z_{11}+Z_{14})}{624},
\label{intdist:eq:q15-third}\\
\zeta(2)
&=\frac{Z_1+Z_2+Z_4+Z_7+Z_8+Z_{11}+Z_{13}+Z_{14}}{192}.
\label{intdist:eq:q15-endpoint}
\end{align}
In the primitive order $(1,2,4,7,8,11,13,14)/15$, the first two
coefficient tuples are respectively
\[
 \frac1{6560}(729,81,9,81,1,729,1,9),\qquad
 \frac1{624}(25,1,25,25,1,1,25,1).
\]
The least common denominator of the full level-fifteen extension
matrix is $511680$, its primitive integral index is $4093440$, and
its nontrivial Smith factors are $2,4,511680$.
\end{corollary}

\begin{proof}
Hurwitz multiplication uses the weights $t_p=p^s$, so here $a=9$ and
$b=25$. The same permutation cycles used in the primitive-cokernel
proof give an elementary elimination derivation of the formulas.
Put $F_j=\zeta(2,j/5)$, and use multiplication by three. The four
relations, in cycle order $1,2,4,3$, are
\[
\begin{aligned}
9F_1&=Z_1+Z_{11}+F_2,&
9F_2&=Z_2+Z_7+F_4,\\
9F_4&=Z_4+Z_{14}+F_3,&
9F_3&=Z_8+Z_{13}+F_1.
\end{aligned}
\]
Eliminating $F_2,F_4,F_3$ gives
\eqref{intdist:eq:q15-fifth}, since $9^4-1=6560$.
Multiplication by five gives the two-cycle
\[
\begin{aligned}
25\zeta(2,1/3)&=Z_1+Z_4+Z_7+Z_{13}+\zeta(2,2/3),\\
25\zeta(2,2/3)&=Z_2+Z_8+Z_{11}+Z_{14}+\zeta(2,1/3).
\end{aligned}
\]
Elimination gives \eqref{intdist:eq:q15-third}, since $25^2-1=624$.
Finally, sorting the absolutely convergent Dirichlet series by
residues prime to fifteen gives
\[
 \sum_{u\in U(15)}\zeta(2,u/15)
 =15^2(1-3^{-2})(1-5^{-2})\zeta(2)=192\zeta(2).
\]
This proves the endpoint formula.

The cokernel parameters are $S_3=26$, $S_5=820$, $D_3=624$,
$D_5=6560$, $g_1=2$, and $g_2=8$. Therefore the three nonunit Smith
factors are $2,4,(624\cdot6560)/8=511680$ and the index is
$624\cdot6560=4093440$. Equivalently, the exact denominator is
\[
 \operatorname{lcm}(624,6560,192)=511680.
\]
The displayed formulas themselves show why every factor in this least
common multiple is necessary: their coefficient lists contain
$1/6560$, $1/624$, and $1/192$ in reduced form. The theorem controls
all the other rows at once. Every zeta value in this example is
finite; no Laurent finite-part convention is needed.
\end{proof}

\subsection{Research directions and integration scope}

The following questions are natural next steps.
\begin{enumerate}
\item Determine the integral primitive-grid cokernel at three or more
prime factors. The permutation-cycle calculation at two primes
suggests a finite incidence presentation, but the resulting torsion
need not be read from separate Euler factors.
\item Extend the explicit cokernel classification to prime powers
$p^a\ell^b$ at ordinary and spectral weights. The quotient basis and
resolution above already apply, but the primitive cokernel contains
additional layers; for example, at $q=12$ and ordinary weights its
nontrivial torsion is $(\mathbb Z/2\mathbb Z)^2$.
\item Give exact common denominators for arbitrary $q$, replacing the
available product bound by an invariant-factor formula or a direct
finite residue-lattice criterion.
\item Study reflection quotients over rings in which two is not a
unit. The complex parity decomposition cannot simply be reduced
modulo two; the integral basis makes the resulting presentation
explicit without presupposing a decomposition.
\item Implement integral jet algebras, including characteristic-$r$
truncations, on this same basis. The underlying quotient is free,
whereas primitive coordinate inclusions can exhibit modular torsion
as well as spectral vanishing.
\end{enumerate}

These results refine the formal algebra of the manuscript. They do
not change its proved complex rank formulas and do not establish
arithmetic independence of polylogarithm, zeta, or Stieltjes values.
The source sections inspected for this part contain no identified
counterexample to their stated complex-algebra theorems. Their
coefficient-ring hypothesis is essential to the character proof;
Theorem~\ref{intdist:thm:basis} supplies a different argument when that
hypothesis is removed.

% Suggested bibliography entries:
% \bibitem{intdist:Kubert1979}
% D. S. Kubert, The universal ordinary distribution, Bulletin de la
% Societe mathematique de France 107 (1979), 179--202,
% https://doi.org/10.24033/bsmf.1891.
% \bibitem{intdist:Ouyang2002}
% Y. Ouyang, On the universal norm distribution, arXiv:math/0210297v2,
% 2002, https://arxiv.org/pdf/math/0210297.
% \bibitem{intdist:Ouyang2001}
% Y. Ouyang, with an appendix by G. W. Anderson, Group cohomology of the
% universal ordinary distribution, J. reine angew. Math. 537 (2001),
% 1--32, https://arxiv.org/abs/math/9911032.
